%实验任务架框架架构
\section{项目架构}
\subsection{安全目标}
根据实验指导书所述场景,可以明确该AI智能体协作安全协议应当满足如下安全目标(包括实验要求和额外添加):
\begin{itemize}
    \item [1.] 可验证身份绑定:所有智能体持有CA颁发的证书,接收方通过CA公钥验证证书真伪,从而验证对方身份,杜绝伪造的智能体接入交互环节。
    \item [2.]数据传输机密性:涉及用户隐私的数据,包括行程信息、订单价格、授权凭证数据等均需要保密,杜绝数据泄露风险.
    \item [3.]授权不扩散:Alice生成的授权凭证包括:总预算、授权权限、授权时效、被授权智能体ID,完整内容由Alice签名认证。凭证内容一旦被篡改,任务随即自动终止,杜绝授权扩散、越权操作。同时授权凭证仅单次会话有效,不可复用。
    \item [4.]执行全程可追溯:FlightAgent完成业务执行后,必须生成带自身私钥签名的审计回执，回执包含：
会话 ID、执行时间、订单详情、实际花费、授权凭证摘要、执行智能体 ID。Alice
汇总所有回执后，可通过对应智能体公钥验证回执真伪，所有交互日志、凭证、
回执可本地留存，支持事后全流程审计追溯，不可抵赖。
    \item [5.] 完整性、时效性与防重放:协议中所有关键消息都必须能够验证是最新的、只用一次的、未被篡改的,即效信息内容的完整性校验、防重放、序列号/时间戳验证、会话绑定。
\end{itemize}

\subsection{流程设计}
根据上述安全要求,我们设计的AI智能体协作航班预定场景协议流程如下:

\begin{figure}[H]
\centering

\begin{tikzpicture}[
    >=Latex,
    node distance=3cm,
    every node/.style={font=\small},
    actor/.style={
        draw,
        rounded corners,
        minimum width=2.5cm,
        minimum height=1cm
    }
]

%------------------------------------------------
% Nodes
%------------------------------------------------

\node[actor] (CA) {CA};

\node[actor,
      below left=2cm and 2cm of CA,
      minimum height=3cm] (Alice) {Alice};

\node[actor,
      below right=2cm and 2cm of CA,
      minimum height=3cm] (FA) {FlightAgent};
\node[actor, below=2cm of Alice] (User) {user};
\node[actor, below=2cm of FA] (Airline) {Airline};

%------------------------------------------------
% CA -> Agents
%------------------------------------------------

\draw[->] (CA.south west) -- node[left] {step0} (Alice.north);
\draw[->] (CA.south east) -- node[right] {step0} (FA.north);

%------------------------------------------------
% Alice <-> FlightAgent
%------------------------------------------------

\draw[<->]
($(Alice.north east)!0.10!(Alice.south east)$)
--
node[above]{step1}
($(FA.north west)!0.10!(FA.south west)$);

\draw[->]
($(Alice.north east)!0.30!(Alice.south east)$)
--
node[above]{step2}
($(FA.north west)!0.30!(FA.south west)$);

\draw[<-]
($(Alice.north east)!0.50!(Alice.south east)$)
--
node[above]{step4}
($(FA.north west)!0.50!(FA.south west)$);

\draw[->]
($(Alice.north east)!0.70!(Alice.south east)$)
--
node[above]{step5}
($(FA.north west)!0.70!(FA.south west)$);

\draw[<-]
($(Alice.north east)!0.90!(Alice.south east)$)
--
node[above]{step7}
($(FA.north west)!0.90!(FA.south west)$);
%------------------------------------------------
% FlightAgent <-> Airline
%------------------------------------------------

\draw[<->]
([xshift=-0.3cm]FA.south)
--
node[left] {step3}
([xshift=-0.3cm]Airline.north);

\draw[<->]
([xshift=0.3cm]FA.south)
--
node[right] {step6}
([xshift=0.3cm]Airline.north);

%------------------------------------------------
% User <-> Alice
%------------------------------------------------

\draw[<-]
([xshift=-0.3cm]Alice.south)
--
node[left] {step0}
([xshift=-0.3cm]User.north);

\draw[->]
([xshift=0.3cm]Alice.south)
--
node[right] {step8}
([xshift=0.3cm]User.north);

\end{tikzpicture}

\caption{Multi-Agent Flight Booking Protocol}
\end{figure}
其中CA表示证书认证机构,Alice为用户智能体,FlightAgent为航班预定智能体,user为用户,Airline为航空公司(也可以是智能体)。上述协议具体步骤如下:
\begin{itemize}
    \item [step0] Alice和FlightAgent在CA处注册得到各自的证书certificate,user向Alice发出预定机票的命令:"预订
    下周(2026.06.08-2026.06.16)青岛—纽约往返机票，总预算不超过 10000 元，单人出行",此后仅需等待Alice与FlightAgent协作交互即可,无需user介入操作。

    \textbf{密码学操作：} CA使用Ed25519为每个Agent签发身份证书，证书绑定Agent身份与公钥：
    \begin{coreconcept}[CA证书签发 — Ed25519]
    \begin{align*}
        (sk_A, pk_A) &\leftarrow \text{Ed25519.KeyGen()} \\
        \text{cert\_body}_A &= \{\text{agent\_id}_A,\; \text{pk}_A,\; \text{permissions},\; \text{not\_before},\; \text{not\_after}\} \\
        \sigma_{CA,A} &= \text{Sign}(sk_{CA},\; \text{canon}(\text{cert\_body}_A)) \qquad \text{(Ed25519, 128-bit安全强度)}\\
        \text{Cert}_A &= (\text{cert\_body}_A,\; \sigma_{CA,A})
    \end{align*}
    \end{coreconcept}
    参数: Ed25519 (Edwards25519曲线), 证书有效期 2026-06-01 $\sim$ 2026-12-31.\\
    代码: \texttt{ca/certificate\_authority.py} — \texttt{issue\_agent()}, \texttt{verify\_certificate()}.

    \item [step1] Alice与FlightAgent之间基于STS协议(认证密钥协商)进行身份验证与ECDH会话密钥协商。

    \textbf{密码学操作：}
    \begin{coreconcept}[Phase 1-2: 认证密钥协商 — X25519 ECDH + Ed25519 + HKDF]
    \textbf{ClientHello (Alice $\to$ FlightAgent):}
    \begin{align*}
        (a, A = g^a) &\leftarrow \text{X25519.KeyGen()} \qquad \text{(临时密钥对)}\\
        \text{body}_A &= \{\text{session\_id},\; \text{agent\_id}_A,\; A_{pub},\; \text{timestamp}\} \\
        \sigma_A &= \text{Sign}(sk_A,\; \text{canon}(\text{body}_A)) \qquad \text{(Ed25519)}
    \end{align*}
    \textbf{ServerHello (FlightAgent $\to$ Alice):}
    \begin{align*}
        \text{Verify}(pk_{CA},\; \text{cert\_body}_A,\; \sigma_{CA,A}) &\to pk_A \qquad \text{(CA证书链验证)}\\
        \text{Verify}(pk_A,\; \text{canon}(\text{body}_A),\; \sigma_A) &\qquad \text{(身份签名验证)}\\
        (b, B = g^b) &\leftarrow \text{X25519.KeyGen()} \\
        \text{shared\_secret} &= a \cdot B = b \cdot A = g^{ab} \qquad \text{(X25519 ECDH)}\\
        \text{HKDF-SHA256}&(\text{salt}=\text{session\_id},\; \text{info}=\text{"a2a-secure..."},\; L=96) \\
        &\to k_{enc}[0:32] \parallel k_{mac}[32:64] \parallel k_{log}[64:96]
    \end{align*}
    \end{coreconcept}
    参数: X25519 (Curve25519), HKDF-SHA256, 派生96字节密钥材料, 三把密钥用途分离(加密/完整性/审计).\\
    代码: \texttt{protocol/handshake.py} — \texttt{HandshakeInitiator}, \texttt{HandshakeResponder}; \texttt{crypto/kdf.py} — \texttt{derive\_keys()}.

    \item [step2] Alice生成授权凭证AuthToken,包括预算、行程、授权信息、授权权限、授权时效、被授权智能体ID等,对这些信息计算杂凑值并签名,将得到的AuthorizationCredential发送给FlightAgent。

    \textbf{密码学操作：}
    \begin{coreconcept}[Phase 3: 授权凭证 — Ed25519签名]
    \begin{align*}
        \text{AuthToken} &= \{\text{issuer},\; \text{delegatee},\; \text{session\_id},\; \text{budget\_cny},\; \text{permissions},\; \text{nonce},\; \text{validity},\; \text{itinerary}\} \\
        \sigma_{Auth} &= \text{Sign}(sk_A,\; \text{canon}(\text{AuthToken})) \qquad \text{(Ed25519)}\\
        \text{Passport} &= \{\text{AuthToken},\; \sigma_{Auth},\; \text{Cert}_A\}
    \end{align*}
    \end{coreconcept}
    参数: budget\_cny=10000, permissions=["flight.search","flight.book"], valid\_minutes=10, nonce=token\_hex(16).\\
    代码: \texttt{protocol/auth\_token.py} — \texttt{make\_authorization()}, \texttt{make\_passport()}.

    \item [step3] FlightAgent接收到授权后先验证签名,如果签名不通过则拒绝服务,直接退出;否则,根据授权凭证的信息执行即机票查询操作,将查询结果,包括具体航班、日期、实际价格,时间戳等信息一并计算杂凑值并签名,将签名和加密后的查询信息发给Alice以询问是否购买预定。

    \textbf{密码学操作 (多层验证链):}
    \begin{coreconcept}[Phase 3: 安全信封解密 + 授权验证 + 查询]
    \textbf{FlightAgent 接收侧:}
    \begin{align*}
        &\text{ReplayWindow}.\text{check}(\text{session\_id},\; \text{seq},\; \text{timestamp}) \qquad \text{(防重放, $\pm30s$窗口)}\\
        &\text{HMAC}_{\text{verify}}(k_{mac},\; \text{envelope\_without\_hmac},\; \text{tag}) \qquad \text{(HMAC-SHA256 完整性)}\\
        &\text{AES-256-GCM}.\text{Decrypt}(k_{enc},\; \text{nonce},\; \text{ct},\; \text{AAD}) \to \text{plaintext} \qquad \text{(AEAD解密)}\\
        &\text{Verify}(pk_{CA},\; \text{cert\_body}_A,\; \sigma_{CA,A}) \to pk_A \qquad \text{(CA证书链)}\\
        &\text{Verify}(pk_A,\; \text{canon}(\text{AuthToken}),\; \sigma_{Auth}) \qquad \text{(授权签名验证)}\\
        &\text{检查: delegatee\_id} \stackrel{?}{=} \text{agent:flight},\; \text{session\_id} \stackrel{?}{=} \text{当前会话},\; \text{budget} \leq 10000,\; \text{permissions} \supseteq \{\text{search},\text{book}\}
    \end{align*}
    \textbf{查询结果加密返回:}
    \begin{align*}
        \text{results} &= \text{BookingService}.\text{search\_flights}(\text{origin},\; \text{dest},\; \text{date}) \\
        \sigma_{Result} &= \text{Sign}(sk_F,\; \text{canon}(\{\text{flights},\; \text{timestamp},\; \text{session\_id}\})) \\
        \text{response} &= \text{AES-256-GCM}.\text{Encrypt}(k_{enc},\; \text{results} \parallel \sigma_{Result})
    \end{align*}
    \end{coreconcept}
    参数: AES-256-GCM (nonce=12B, key=32B), HMAC-SHA256 (key=32B), ReplayWindow (skew=30s).\\
    代码: \texttt{protocol/envelope.py} — \texttt{open\_envelope()}; \texttt{agents/flight\_agent.py} — \texttt{process\_secure\_request()}, \texttt{\_handle\_query()}.

    \item [step4] Alice收到FlightAgent发送来的信息后先验证签名,如果签名不通过则拒绝服务并直接退出;否则,根据查询结果,判断是否购买预定,发送BookRequest给FlightAgent允许其预定机票。

    \textbf{密码学操作：}
    \begin{coreconcept}[Phase 4: 查询结果验证 + 预订确认]
    \begin{align*}
        &\text{AES-256-GCM}.\text{Decrypt}(k_{enc},\; \text{response}) \to \text{results} \parallel \sigma_{Result} \\
        &\text{Verify}(pk_{CA},\; \text{Cert}_F) \to pk_F,\qquad \text{Verify}(pk_F,\; \text{canon}(\text{results}),\; \sigma_{Result}) \qquad \text{(查询结果验签)}\\
        &\text{BookRequest} = \{\text{flight\_no},\; \text{confirm}\},\qquad \text{Enc}_{k_{enc}}(\text{BookRequest}) \to \text{FlightAgent}
    \end{align*}
    \end{coreconcept}
    代码: \texttt{agents/alice\_agent.py} — \texttt{query\_flights()}, \texttt{confirm\_booking()}.

    \item [step5] FlightAgent收到Alice发送来的BookRequest信息,同样验证签名,不通过则拒绝服务并退出,否则执行step6。

    \textbf{密码学操作：} 重复step3的安全信封解密与HMAC验证流程(同 \texttt{open\_envelope()}),验签通过后执行预订.

    \item [step6] FlightAgent在航空公司Airline处预定机票,将具体订单ID、航空公司信息、机票信息、具体花费等信息作为收据Receipt,签名并发回给Alice。

    \textbf{密码学操作：}
    \begin{coreconcept}[Phase 5-6: 订单执行 + 签名审计回执]
    \begin{align*}
        \text{order} &= \text{BookingService}.\text{create\_order}(\text{flight\_no},\; \text{passengers},\; \text{route}) \\
        \text{receipt\_body} &= \{\text{session\_id},\; \text{executed\_at},\; \text{order},\; \text{actual\_spend},\; \text{auth\_digest},\; \text{executor\_id}\} \\
        \text{auth\_digest} &= \text{SHA-256}(\text{canon}(\text{AuthToken})) \qquad \text{(授权摘要绑定)}\\
        \sigma_{Receipt} &= \text{Sign}(sk_F,\; \text{canon}(\text{receipt\_body})) \qquad \text{(Ed25519 不可抵赖)}\\
        \text{Klog\_HMAC} &= \text{HMAC}(k_{log},\; \text{canon}(\text{receipt\_body})) \qquad \text{(审计链绑定)}
    \end{align*}
    \end{coreconcept}
    参数: Ed25519签名, SHA-256杂凑, HMAC-SHA256 (key=$k_{log}$, 32B).\\
    代码: \texttt{protocol/receipt.py} — \texttt{make\_receipt()}; \texttt{agents/booking\_service.py} — \texttt{create\_order()}.

    \item [step7] Alice验证签名,保存审计记录,并将结果报告给user。

    \textbf{密码学操作：}
    \begin{coreconcept}[Phase 6: Alice侧回执验证 — 双重验证]
    \begin{align*}
        &\text{Verify}(pk_{CA},\; \text{Cert}_F) \to pk_F \qquad \text{(CA证书链验证)}\\
        &\text{Verify}(pk_F,\; \text{canon}(\text{receipt\_body}),\; \sigma_{Receipt}) \qquad \text{(Ed25519 回执验签)}\\
        &\text{HMAC}_{\text{verify}}(k_{log},\; \text{canon}(\text{receipt\_body}),\; \text{Klog\_HMAC}) \qquad \text{(审计链完整性)}
    \end{align*}
    通过后写入审计日志: \texttt{AuditLogger.log(event\_type="receipt", detail\_digest=SHA-256(receipt), klog\_hmac)}.
    \end{coreconcept}
    代码: \texttt{agents/alice\_agent.py} — \texttt{verify\_receipt()}; \texttt{audit/logger.py} — \texttt{AuditLogger}; \texttt{audit/store.py} — \texttt{AuditStore}.

    \item [step8] Alice将最终结果报告给用户,包括成功预定或中途发生的异常信息。
\end{itemize}

\subsection{密码学算法总览}
协议使用的密码学原语及参数汇总:
\begin{table}[H]
\centering
\begin{tabular}{@{}llll@{}}
\toprule
\textbf{用途} & \textbf{算法} & \textbf{参数} & \textbf{安全强度} \\
\midrule
身份签名 & Ed25519 (Edwards25519) & 256-bit密钥, 512-bit签名 & 128-bit \\
密钥协商 & X25519 ECDH & 256-bit密钥, Curve25519 & 128-bit \\
密钥派生 & HKDF-SHA256 & salt=session\_id, L=96B & 128-bit \\
数据加密 & AES-256-GCM & 256-bit密钥, 96-bit nonce & 256-bit \\
完整性校验 & HMAC-SHA256 & 256-bit密钥 & 256-bit \\
消息摘要 & SHA-256 & 256-bit输出 & 128-bit (抗碰撞) \\
防重放 & ReplayWindow & $\pm30s$时钟偏移窗口, (sid,seq)去重 & — \\
审计绑定 & Klog HMAC-SHA256 & 密钥$k_{log}$(32B), 绑定回执与日志 & 256-bit \\
\bottomrule
\end{tabular}
\caption{协议密码学算法总览}
\end{table}
